\section{Appendix}\label{sec:APPENDIX}
\subsection*{Introduction}
Due to its nature as a noble gas, \(^{222}\text{Rn}\) is difficult to remove from the detector environment, making it a persistent background source also in detector components. Among its progeny, several isotopes are particularly important:

\begin{itemize}
  \item \textbf{Bismuth (\(^{214}\text{Bi}\))}: A short--lived beta emitter with a high Q--value (up to \SI{3.27}{\mega\electronvolt}) and multiple gamma lines. It is the most significant source of high--energy background from \(^{222}\text{Rn}\)~\cite{pdg2024}.
  \item \textbf{Bismuth (\(^{210}\text{Bi}\))}: A beta emitter that produces a continuum low--energy background up to \SI{1.16}{\mega\electronvolt}.
  \item \textbf{Lead (\(^{210}\text{Pb}\))}: A long--lived isotope (\(t_{1/2} = 22.3\)~years) that accumulates on detector surfaces~\cite{exo_radon}.
  \item \textbf{Polonium (\(^{210}\text{Po}\))}: An alpha emitter with a half--life of 138~days. It can be used as a surface contamination proxy due to its distinct alpha signal~\cite{saldanha_thesis}.
\end{itemize}